\documentclass[10pt,a4paper]{amsart}
\usepackage[margin=19mm]{geometry}
\usepackage{amsmath,amssymb,mathtools}
\usepackage[scr=rsfs,cal=euler]{mathalfa}
\usepackage{xfrac}
\usepackage[unicode,hidelinks,bookmarks=false]{hyperref}
\newcommand{\ie}{i.e.\ }
\newcommand{\cf}{cf.\ }
\newcommand{\resp}{respectively\ }
\newcommand{\Subsection}{\subsection}
\providecommand{\nobreakdash}{\nobreak\hbox{-}}
\setlength{\emergencystretch}{2em}
\multlinegap0pt
\usepackage{bm}
\usepackage{bbm}
\usepackage{dsfont}
\usepackage{xspace}
\usepackage{cancel}
\usepackage{mathtools}
\usepackage{amsthm}
\makeatletter
\@ifundefined{theorem}{\newtheorem{theorem}{Theorem}}{}
\@ifundefined{lemma}{\newtheorem{lemma}[theorem]{Lemma}}{}
\@ifundefined{proposition}{\newtheorem{proposition}[theorem]{Proposition}}{}
\makeatother
\newcommand{\Ahat}{\widehat{A}_G}
\newcommand{\Cay}{\mathrm{Cay}}
\newcommand{\Shat}{\hat{S}^{(j)}}
\newcommand{\Uj}{U_G^{(j)}}
\newcommand{\Z}{\mathbb{Z}}
\newcommand{\braket}[2]{\langle #1 | #2 \rangle}
\newcommand{\bra}[1]{\langle #1|}
\newcommand{\eqdef}{\coloneqq}
\newcommand{\ket}[1]{|#1\rangle}
\newcommand{\phij}{\phi_j}
\newcommand{\srg}{\mathrm{srg}}
\newcommand{\vplus}{v_+}
\title[The Quantum Walk Characteristic Polynomial Distinguishes All]{The Quantum Walk Characteristic Polynomial Distinguishes All Strongly Regular Graphs of Prime Orde}
\author{Diego Roldan}
\date{Source: arXiv:2604.01507v1 (CC BY 4.0)}
\begin{document}
\maketitle

\noindent\emph{Source and licence.} The Quantum Walk Characteristic Polynomial Distinguishes All Strongly Regular Graphs of Prime Orde. Version arXiv:2604.01507v1 (CC BY 4.0), distributed under the Creative Commons Attribution 4.0 International licence, as recorded in the arXiv licence metadata for this exact version and shown on the abstract page https://arxiv.org/abs/2604.01507v1. Full rights notice: licenses/arxiv-2604.01507-CC-BY-4.0.txt.\par\smallskip

\noindent\textit{Editorial scope.} Counted unit: the Lemma whose source label is \texttt{lem:formula} in the article (source0.tex lines 370--380, source label \texttt{lem:formula}), delivered together with the article's own complete original proof of it (source0.tex lines 382--437, one \texttt{proof} environment of 56 lines). No mathematics is written, repaired or re-derived: statement and proof are the article's own text line for line. Required context retained uncounted: eq:Shat (lines 257--259); prop:key (lines 292--299); lem:lindep (lines 327--338). Every definition, display and label of those lines is retained unchanged. Omitted from this package and not redistributed: 1--256, 260--291, 300--326, 339--369, 381--381, 438--700. Nothing inside a retained excerpt is omitted or altered: every retained body is delivered exactly as the article's lines are written, so no omission receipt of any kind accompanies this package. Citations: the retained excerpts carry no citation command at all, so no bibliography entry of the article is transcribed and no reference is redistributed. Reference adaptation: none was needed and none was made; every reference inside the delivered unit resolves inside it, so no unresolved reference or citation is produced. Printed slips kept literal: no printed word, symbol, hypothesis or number of the counted statement or of its proof was altered. Layout and class: the article's own class is \texttt{article} with packages including amsmath, amssymb, amsthm, hyperref, graphicx, mathtools, mathrsfs, booktabs, enumerate, geometry; the wrapper is the lane's pinned \texttt{amsart} editorial wrapper at 10pt on A4, which restates the theorem-like environments the retained text needs and adds the article's own macro definitions that the retained text needs (\texttt{\textbackslash Ahat}, \texttt{\textbackslash Cay}, \texttt{\textbackslash Shat}, \texttt{\textbackslash Uj}, \texttt{\textbackslash Z}, \texttt{\textbackslash bra}, \texttt{\textbackslash braket}, \texttt{\textbackslash eqdef}, \texttt{\textbackslash ket}, \texttt{\textbackslash phij}, \texttt{\textbackslash srg}, \texttt{\textbackslash vplus}), with extra wrapper packages (bm, bbm, dsfont, xspace, cancel, mathtools). The delivered numbering is the unit's own. Assets: no retained span contains a figure environment, a graphics-inclusion command, a PDF-page inclusion, a table or a TikZ picture; no image file is copied into this package, the package declares no asset dependency and no image file is redistributed with it. Licence: version arXiv:2604.01507v1 of this article is distributed under the Creative Commons Attribution 4.0 International licence, as recorded in the arXiv licence metadata for this exact version and shown on the abstract page https://arxiv.org/abs/2604.01507v1. A full rights notice is delivered at licenses/arxiv-2604.01507-CC-BY-4.0.txt. Characters: every retained excerpt is pure ASCII, so the delivered engine, XeLaTeX with its default Latin Modern text font, reports no missing character.
  \begin{equation}\label{eq:Shat}
    \Shat\ket{s} = \omega^{js}\ket{-s}, \qquad s \in S.
  \end{equation}

\begin{proposition}[Key Identity]\label{prop:key}
  For every $j \in \Z_p$, let
  $\ket{\phij} \eqdef \frac{1}{\sqrt{k}}\sum_{s \in S}\omega^{js}\ket{s}$.
  Then
  \[
    \bra{\phij}\Shat\ket{\phij} = \frac{\Ahat(j)}{k}.
  \]
\end{proposition}

\begin{lemma}[Linear Independence]\label{lem:lindep}
  Let $G = \Cay(\Z_p, S)$ be an $\srg(p,k,\lambda,\mu)$ with $p$ prime
  and $k \geq 2$.  For every $j \in \Z_p \setminus \{0\}$, the vectors
  \[
    \ket{\phij} = \frac{1}{\sqrt{k}}\sum_{s \in S}\omega^{js}\ket{s}
    \qquad\text{and}\qquad
    \Shat\ket{\vplus} = \frac{1}{\sqrt{k}}\sum_{s \in S}\omega^{js}\ket{-s}
  \]
  are linearly independent, so
  $W_j \eqdef \mathrm{span}\{\ket{\phij},\Shat\ket{\vplus}\}$
  has dimension exactly $2$.
\end{lemma}

\begin{lemma}[Explicit Formula]\label{lem:formula}
  For every $j \in \Z_p$ with $\Ahat(j)/k \neq \pm 1$:
  \begin{equation}\label{eq:chi-block}
    \chi_q\!\bigl(\Uj, \lambda\bigr)
    = (\lambda - 1)^{(k-2)/2}(\lambda + 1)^{(k-2)/2}
      \!\left(\lambda^2 - \frac{2\Ahat(j)}{k}\,\lambda + 1\right).
  \end{equation}
  In particular, the eigenvalues of $\Uj$ are
  $\bigl\{+1^{(k-2)/2},\,-1^{(k-2)/2},\,e^{i\theta_j},\,e^{-i\theta_j}\bigr\}$
  where $\cos\theta_j = \Ahat(j)/k$.
\end{lemma}

\begin{proof}
  Set $c_j \eqdef \Ahat(j)/k \in (-1,1)$.

  \medskip
  \noindent\textbf{Invariant subspace.}
  By Lemma~\ref{lem:lindep}, $W_j = \mathrm{span}\{\ket{\phij},
  \Shat\ket{\vplus}\}$ has dimension $2$.  Write
  $\ket{\phij} = c_j\ket{\vplus}+\ket{\phij^\perp}$ where
  $\braket{\vplus}{\phij}=c_j$ (Proposition~\ref{prop:key}).  Then
  \[
    \Uj\ket{\phij}
    = \Shat(c_j\ket{\vplus}-\ket{\phij^\perp})
    = c_j\,\Shat\ket{\vplus} - \Shat\ket{\phij^\perp}.
  \]
  Since $\Shat$ is unitary, $\ket{\phij^\perp}\perp\ket{\phij}$, and
  $\ket{\phij^\perp}\perp\ket{\vplus}$, the vector $\Shat\ket{\phij^\perp}$
  lies in $W_j^\perp$.  So $\Uj\ket{\phij}\in W_j$.  A symmetric argument
  gives $\Uj(\Shat\ket{\vplus})\in W_j$, so $W_j$ is $\Uj$-invariant.

  \medskip
  \noindent\textbf{Characteristic polynomial on $W_j$.}
  The matrix of $\Uj|_{W_j}$ in the Gram--Schmidt orthonormal basis of
  $W_j$ has trace $2c_j$ (Proposition~\ref{prop:key}) and determinant $1$
  (unitarity), giving characteristic polynomial
  $\lambda^2-2c_j\lambda+1$ with roots $e^{\pm i\theta_j}$.

  \medskip
  \noindent\textbf{Eigenvalues on $W_j^\perp$.}
  We compute $\operatorname{tr}(\Uj)$ directly.  From
  $\Uj = \Shat C_{\mathrm{loc}}$ and equation~\eqref{eq:Shat}, the
  diagonal entry at $\ket{s}$ is
  \[
    (\Uj)_{s,s} = \omega^{js}(C_{\mathrm{loc}})_{-s,\,s}
                = \omega^{js}\cdot\frac{2}{k},
  \]
  since $s \neq -s$ for all $s \in S \subset \Z_p^*$.
  Summing over $S$:
  \[
    \operatorname{tr}(\Uj) = \frac{2}{k}\sum_{s\in S}\omega^{js} = 2c_j.
  \]
  The restriction $\Uj|_{W_j}$ has trace $2c_j$ (from
  $\lambda^2-2c_j\lambda+1$), so
  \[
    \operatorname{tr}\!\bigl(\Uj|_{W_j^\perp}\bigr) = 2c_j - 2c_j = 0.
  \]
  Every $\ket{w}\in W_j^\perp$ satisfies $\braket{\vplus}{w}=0$, hence
  $C_{\mathrm{loc}}\ket{w}=-\ket{w}$, and therefore
  $\Uj|_{W_j^\perp} = -\Shat|_{W_j^\perp}$.
  Since $(\Shat)^2 = I_k$, the eigenvalues of $-\Shat|_{W_j^\perp}$ lie
  in $\{+1,-1\}$.  Their sum is zero and
  $\dim W_j^\perp = k-2$, so there are exactly $\tfrac{k-2}{2}$
  eigenvalues $+1$ and $\tfrac{k-2}{2}$ eigenvalues $-1$.
  Combined with $e^{\pm i\theta_j}$ from $W_j$, the full spectrum is
  $\bigl\{+1^{(k-2)/2},\,-1^{(k-2)/2},\,e^{i\theta_j},\,e^{-i\theta_j}
  \bigr\}$, which gives~\eqref{eq:chi-block}.
\end{proof}

\end{document}
